\section*{Chapitre \ref{ch:b2:angiosperm-reproduction} --- La reproduction sexuée des plantes à fleurs}

\begin{solution}{exo:b2:angiosperm-reproduction:1}
Sépales (calice), pétales (corolle), \omterm{def:b2:angiosperm-reproduction:flower}{étamines} (androcée), \omterm{def:b2:angiosperm-reproduction:flower}{carpelles}
(gynécée). \omterm{def:b2:angiosperm-reproduction:flower}{Étamine}~: filet et anthère. \omterm{def:b2:angiosperm-reproduction:flower}{Carpelle}~: \omterm{def:b2:angiosperm-reproduction:flower}{stigmate}, style,
\omterm{def:b2:angiosperm-reproduction:flower}{ovaire} contenant les \omterm{def:b2:plant-life-cycles:heterospory}{ovules}. Toutes ces pièces appartiennent au
\omterm{def:b2:plant-life-cycles:cycles}{sporophyte} \omterm{def:b2:meiosis-heredity:meiosis}{diploïde}~; seuls le \omterm{prop:b2:angiosperm-reproduction:pollen}{grain de pollen} (dans l'anthère) et le
\omterm{prop:b2:angiosperm-reproduction:embryosac}{sac embryonnaire} (dans l'\omterm{def:b2:plant-life-cycles:heterospory}{ovule}) sont des gamétophytes \omterm{def:b2:meiosis-heredity:meiosis}{haploïdes}.
\end{solution}

\begin{solution}{exo:b2:angiosperm-reproduction:2}
\omterm{def:b2:plant-life-cycles:heterospory}{Microspore} $n$~; cellule végétative $n$~; cellule spermatique $n$~;
oosphère $n$~; noyau polaire $n$~; zygote $2n$~; albumen $3n$~;
\omterm{def:b2:angiosperm-reproduction:seed}{tégument de la graine} $2n$ (téguments maternels)~; paroi du \omterm{def:b2:angiosperm-reproduction:seed}{fruit} $2n$
(\omterm{def:b2:angiosperm-reproduction:flower}{ovaire} maternel).
\end{solution}

\begin{solution}{exo:b2:angiosperm-reproduction:3}
Vent (graminées, chêne, bouleau, noisetier, plantain)~: petites \omterm{def:b2:angiosperm-reproduction:flower}{fleurs}
vertes, sans pétales, sans parfum ni \omterm{def:b2:angiosperm-reproduction:pollination}{nectar}, anthères pendant au bout
de longs filets, \omterm{def:b2:angiosperm-reproduction:flower}{stigmates} plumeux, énormes quantités de \omterm{def:b2:plant-life-cycles:heterospory}{pollen} lisse
et sec, floraison avant l'apparition des feuilles. Abeille (trèfle,
sauge, digitale, lavande, pommier)~: pétales colorés, souvent avec des
guides ultraviolets, piste d'atterrissage, parfum, \omterm{def:b2:angiosperm-reproduction:pollination}{nectar}, \omterm{def:b2:plant-life-cycles:heterospory}{pollen}
collant et ornementé en quantité modérée, floraison en plein jour.
\end{solution}

\begin{solution}{exo:b2:angiosperm-reproduction:4}
Les deux cellules spermatiques d'un même \omterm{thm:b2:angiosperm-reproduction:tube}{tube pollinique} fusionnent,
l'une avec l'oosphère (zygote, $2n$, l'embryon), l'autre avec la cellule
centrale (noyau de l'albumen, $3n$, la réserve nutritive). L'albumen ne
commence à se développer qu'après la fécondation~: la plante
n'approvisionne donc que les \omterm{def:b2:plant-life-cycles:heterospory}{ovules} qui contiennent un embryon,
contrairement aux gymnospermes, qui construisent leur réserve avant.
\end{solution}

\begin{solution}{exo:b2:angiosperm-reproduction:5}
$25/1.2 = \qty{21}{h}$. Les soies émergeant les jours 5, 6 et 7
rencontrent du \omterm{def:b2:plant-life-cycles:heterospory}{pollen}~; celles des jours 8 à 12 non~: $3/8$ des soies
sont pollinisées, et l'épi est stérile aux cinq huitièmes.
\end{solution}

\begin{solution}{exo:b2:angiosperm-reproduction:6}
$S_1S_2\times S_1S_2$~: 0. $S_1S_2\times S_1S_3$~: $\tfrac12$ (seul le
\omterm{def:b2:plant-life-cycles:heterospory}{pollen} $S_3$ croît). $S_1S_2\times S_3S_4$~: tout. Descendants du
deuxième croisement~: oosphères $S_1$ ou $S_2$, cellules spermatiques
$S_3$~: $S_1S_3$ et $S_2S_3$, moitié chacun.
\end{solution}

\begin{solution}{exo:b2:angiosperm-reproduction:7}
Une plante porte 2 des $n$ \omterm{def:b2:meiosis-heredity:vocabulary}{allèles} et rejette le \omterm{def:b2:plant-life-cycles:heterospory}{pollen} porteur de l'un
ou de l'autre~: fraction compatible $(n - 2)/n$. Un \omterm{def:b2:meiosis-heredity:vocabulary}{allèle} nouveau
n'est rejeté par aucun pistil~: son \omterm{def:b2:plant-life-cycles:heterospory}{pollen} engendre plus de \omterm{def:b2:angiosperm-reproduction:seed}{graines} que
la moyenne, et il se répand jusqu'à être aussi fréquent que les autres~;
le même avantage protège tout \omterm{def:b2:meiosis-heredity:vocabulary}{allèle} rare de la disparition. La
sélection accumule donc les \omterm{def:b2:meiosis-heredity:vocabulary}{allèles}, et les populations naturelles
d'espèces auto-incompatibles en portent des dizaines.
\end{solution}

\begin{solution}{exo:b2:angiosperm-reproduction:8}
Les deux \omterm{prop:b2:angiosperm-reproduction:embryosac}{noyaux polaires} sont des copies d'une même \omterm{def:b2:plant-life-cycles:heterospory}{macrospore}~: la
cellule centrale est donc $YY$ ou $yy$ (moitié chacune)~; la cellule
spermatique est $Y$ ou $y$ (moitié chacune)~: albumen $YYY$, $YYy$,
$Yyy$, $yyy$, à $\tfrac14$ chacun~; trois quarts de grains jaunes.
Pollinisée par $yy$~: $YYy$ et $yyy$, moitié chacun --- la moitié des
grains sont jaunes.
\end{solution}

\begin{solution}{exo:b2:angiosperm-reproduction:9}
Des grains d'orge sont coupés en deux~; les moitiés portant l'embryon
digèrent leur amidon, les moitiés sans embryon non, sauf si l'on ajoute
de la gibbérelline, auquel cas elles le digèrent. (a) L'embryon est la
source du signal, non de l'enzyme. (b) La gibbérelline est le signal,
suffisant à lui seul. (c) L'amylase est fabriquée par la couche à
aleurone de l'albumen, qui répond à l'hormone. La moitié sans embryon
sépare le signal de la réponse~: sans elle, on ne pourrait pas savoir si
l'embryon digère lui-même l'amidon.
\end{solution}

\begin{solution}{exo:b2:angiosperm-reproduction:10}
Vent~: $10^{6}\times\qty{1}{\micro g} = \qty{1}{g}$ par \omterm{def:b2:plant-life-cycles:heterospory}{ovule}. Animal~:
\qty{1}{mg} de \omterm{def:b2:plant-life-cycles:heterospory}{pollen} plus \qty{0.5}{mg} de sucre $= \qty{1.5}{mg}$
par \omterm{def:b2:plant-life-cycles:heterospory}{ovule} --- environ sept cents fois moins cher. La \omterm{def:b2:angiosperm-reproduction:pollination}{pollinisation} par
le vent n'a besoin d'aucun partenaire~: elle fonctionne au début du
printemps, avant que les insectes ne volent, dans les milieux froids,
venteux et ouverts, la nuit et sous la pluie, dans les peuplements
denses d'une seule espèce, et elle ne peut être ni volée par des
pilleurs de \omterm{def:b2:angiosperm-reproduction:pollination}{nectar} ni abandonnée par un pollinisateur qui disparaît.
\end{solution}

\begin{solution}{exo:b2:angiosperm-reproduction:11}
Les \omterm{def:b2:angiosperm-reproduction:seed}{graines} qui tombent au pied de l'arbre parent atterrissent dans son
ombre, parmi ses racines, là où se concentrent ses prédateurs de \omterm{def:b2:angiosperm-reproduction:seed}{graines}
et ses pathogènes spécifiques~; presque toutes meurent. Un oiseau
emporte une \omterm{def:b2:angiosperm-reproduction:seed}{graine} à des kilomètres (\qty{12}{km} en vingt minutes) et
la dépose, avec de l'engrais, dans une haie ou une clairière~; même si
une \omterm{def:b2:angiosperm-reproduction:seed}{graine} sur cent seulement tombe en un bon site, c'est plus que
toutes les \omterm{def:b2:angiosperm-reproduction:seed}{graines} tombées sous l'arbre, et la population s'étend à la
vitesse des oiseaux, non à celle des \omterm{def:b2:angiosperm-reproduction:seed}{fruits} qui tombent. La chair est le
prix du billet.
\end{solution}

\begin{solution}{exo:b2:angiosperm-reproduction:12}
L'enfermement oblige le \omterm{def:b2:plant-life-cycles:heterospory}{pollen} à croître à travers des tissus maternels,
ce qui permet au pistil de le tester (auto-incompatibilité, rejet des
espèces étrangères), de laisser de nombreux tubes entrer en compétition
pour que le plus rapide engendre la \omterm{def:b2:angiosperm-reproduction:seed}{graine}, de protéger les \omterm{def:b2:plant-life-cycles:heterospory}{ovules} de la
dessiccation et des herbivores, et de transformer l'\omterm{def:b2:angiosperm-reproduction:flower}{ovaire} en un \omterm{def:b2:angiosperm-reproduction:seed}{fruit}
qui recrute des animaux pour la dispersion~; la \omterm{prop:b2:angiosperm-reproduction:double}{double fécondation}
n'approvisionne ensuite que les \omterm{def:b2:plant-life-cycles:heterospory}{ovules} fécondés. Les coûts~: un \omterm{prop:b2:angiosperm-reproduction:pollen}{grain
de pollen} doit atteindre un \omterm{def:b2:angiosperm-reproduction:flower}{stigmate} plutôt que l'\omterm{def:b2:plant-life-cycles:heterospory}{ovule} lui-même, ce qui
exige un style, un tube et son guidage, et le \omterm{def:b2:angiosperm-reproduction:seed}{fruit} est un
investissement lourd. Le bilan a si fortement favorisé les angiospermes
qu'elles sont passées de rien aux neuf dixièmes des espèces végétales en
cent millions d'années.
\end{solution}

\begin{solution}{pb:b2:angiosperm-reproduction:1}
\textbf{1.} Sur A ($S_1S_2$)~: \omterm{def:b2:plant-life-cycles:heterospory}{pollen} de A $f = 0$~; de B ($S_1$, $S_3$)
$f =
\tfrac12$~; de C ($S_3$, $S_4$) $f = 1$.
\textbf{2.} Sur B ($S_1S_3$)~: A $\tfrac12$, B 0, C $\tfrac12$. Sur C
($S_3S_4$)~: A 1, B $\tfrac12$, C 0.
\textbf{3.} 50 grains compatibles~; $0.1\times 50 = 5$ \omterm{def:b2:plant-life-cycles:heterospory}{ovules}.
\textbf{4.} Depuis C~: 100 compatibles, $0.1\times 100 = 10$,
c'est-à-dire les dix \omterm{def:b2:plant-life-cycles:heterospory}{ovules}~; depuis A~: aucun.
\textbf{5.} Les \omterm{def:b2:angiosperm-reproduction:seed}{fruits} issus de visites depuis C (10 pépins) sont
conservés~; depuis B (5 pépins), tout juste conservés~; depuis A, ils
tombent. Un verger de la seule variété A ne produit rien.
\textbf{6.} La moitié des \omterm{def:b2:angiosperm-reproduction:flower}{fleurs} nouent un \omterm{def:b2:angiosperm-reproduction:seed}{fruit} (celles visitées depuis
C)~: 2500 \omterm{def:b2:angiosperm-reproduction:seed}{fruits} par arbre, dix fois les 250 nécessaires --- une pleine
récolte, et l'arbre fera tomber le surplus.
\textbf{7.} Il libère du \omterm{def:b2:plant-life-cycles:heterospory}{pollen} compatible pendant toute la floraison de
chaque variété~; mais s'il partage ses deux \omterm{def:b2:meiosis-heredity:vocabulary}{allèles} $S$ avec une
variété, il ne pollinise rien sur elle.
\textbf{8.} $8\times 10^{7}/(7\times 86400) = 132$ grains par mètre carré
et par seconde.
\textbf{9.} $132/0.25 = 530$ grains par mètre cube.
\textbf{10.} $132\times 10^{-4} = 0.013$ par seconde~: un grain toutes
les \qty{76}{s}.
\textbf{11.} $0.013\times 86400 \approx 1100$ grains par jour~; le
premier tube qui atteint l'\omterm{def:b2:plant-life-cycles:heterospory}{ovule} le féconde, et l'\omterm{def:b2:plant-life-cycles:heterospory}{ovule} est alors fermé
aux autres.
\textbf{12.} \qty{20}{h} après la \omterm{def:b2:angiosperm-reproduction:pollination}{pollinisation}.
\textbf{13.} \omterm{def:b2:plant-life-cycles:heterospory}{Pollen} les jours 0--7, soies à partir du jour 5~: seuls les
jours 5, 6 et 7 se recouvrent --- $3/8$ des soies, et
les soies qui émergent après le jour 7 ne reçoivent rien~; l'épi est en
grande partie stérile. La sécheresse à la floraison femelle est la cause
classique de l'échec d'une récolte de maïs.
\textbf{14.} Embryon $Yy$~; albumen $YYy$.
\textbf{15.} Jaune, car $Y$ est \omterm{def:b2:meiosis-heredity:vocabulary}{dominant}~: le caractère du parent
pollinisateur se manifeste dans la \omterm{def:b2:angiosperm-reproduction:seed}{graine} portée par la plante mère ---
c'est la xénie.
\textbf{16.} Cellule centrale $YY$ ou $yy$~: albumen $YYy$ (moitié) et
$yyy$ (moitié).
\textbf{17.} $YYy$~: 20 unités~; $Yyy$~: 10 unités~; un $yyY$ dont le $Y$
vient du père a le même \omterm{def:b2:meiosis-heredity:vocabulary}{génotype} $Yyy$, 10 unités. Deux doses
paternelles ne peuvent pas exister~: une seule cellule spermatique
féconde la cellule centrale, et l'albumen compte toujours deux génomes
maternels pour un paternel.
\textbf{18.} C'est la réserve d'amidon~; $0.3/0.33 = \qty{91}{\%}$ du
grain.
\textbf{19.} La réserve n'est construite qu'après la fécondation~: aucune
ressource n'est investie dans des \omterm{def:b2:plant-life-cycles:heterospory}{ovules} sans embryon~; le gamétophyte
des gymnospermes est construit à l'avance et perdu si la \omterm{def:b2:angiosperm-reproduction:pollination}{pollinisation}
échoue.
\textbf{20.} $0.9\times 600 = 540$ grains par épi et par plante~;
$4320$ par mètre carré~; $4320\times 0.3 = \qty{1.3}{kg}$ d'albumen
par mètre carré.
\textbf{21.} \qty{13}{t/ha}.
\textbf{22.} Carbone du grain $0.5\times 1.2 = \qty{0.6}{kg/m^2}$~;
carbone de l'albumen $0.45\times 1.3 = \qty{0.58}{kg/m^2}$~: cohérent
(le petit reste correspond à l'embryon et au tégument).
\textbf{23.} $8\times 10^{7}/4320 \approx \num{18500}$ \omterm{prop:b2:angiosperm-reproduction:pollen}{grains de pollen}
par grain récolté.
\textbf{24.} Son propre nuage de \omterm{def:b2:plant-life-cycles:heterospory}{pollen} est ténu et emporté par le vent,
et sa panicule libère surtout son \omterm{def:b2:plant-life-cycles:heterospory}{pollen} avant l'émergence de ses propres
soies~: de nombreuses soies ne sont jamais pollinisées et l'épi présente
des manques.
\textbf{25.} Bloc A~: la moitié des \omterm{def:b2:angiosperm-reproduction:flower}{fleurs} nouent un \omterm{def:b2:angiosperm-reproduction:seed}{fruit}, 2500 par
arbre, une pleine récolte. Champ~: 4320 grains par mètre carré,
\qty{13}{t/ha}.
\end{solution}
